\section{Вещественные аналитические функции}

\begin{definition}
Функция $f\colon I\to\R$ ($I$ — интервал) называется \emph{вещественно аналитической} в точке
$x_0\in I$, если существует окрестность $(x_0-R,x_0+R)$, на которой $f$ представляется сходящимся
степенным рядом
\[
  f(x)=\sum_{n=0}^\infty a_n(x-x_0)^n .
\]
$f$ аналитична на $I$, если она аналитична в каждой его точке. Класс таких функций обозначают
$C^{\omega}$.
\end{definition}

\begin{example}
Многочлен $P_n(x)=\sum_{k=0}^n a_k(x-x_0)^k$; функция $\dfrac1{1-x}=\sum_{n\ge0}x^n$ аналитична в
нуле ($\abs x<1$); $e^x,\ \sin x,\ \cos x$ аналитичны при всех $x_0\in\R$.
\end{example}

\section{Дифференцирование. Коэффициенты Тейлора}

\begin{theorem}
Если $f$ аналитична в $x_0$ и $f(x)=\sum a_n(x-x_0)^n$ с радиусом сходимости $R$, то на
$(x_0-R,x_0+R)$:
\begin{enumerate}[label=\arabic*),leftmargin=*]
  \item $f$ бесконечно дифференцируема ($f\in C^\infty$);
  \item $\displaystyle f^{(k)}(x)=\sum_{n=k}^\infty a_n\,n(n-1)\cdots(n-k+1)(x-x_0)^{n-k}$;
  \item ряд для $f^{(k)}$ имеет тот же радиус сходимости.
\end{enumerate}
\end{theorem}

\begin{theorem}[коэффициенты Тейлора]
Коэффициенты аналитической функции однозначно определяются её производными:
\[
  a_k=\frac{f^{(k)}(x_0)}{k!},\qquad\text{поэтому}\qquad
  f(x)=\sum_{k=0}^\infty\frac{f^{(k)}(x_0)}{k!}(x-x_0)^k .
\]
\end{theorem}

\begin{proof}
Подставим $x=x_0$ в формулу для $f^{(k)}$ из предыдущей теоремы: все слагаемые с $n>k$ содержат
$(x-x_0)^{n-k}$ и обращаются в нуль, остаётся член с $n=k$:
$f^{(k)}(x_0)=a_k\cdot k!$, откуда $a_k=f^{(k)}(x_0)/k!$.
\end{proof}

\begin{definition}
Для $f\in C^\infty$ ряд $\displaystyle\sum_{k=0}^\infty\frac{f^{(k)}(x_0)}{k!}(x-x_0)^k$ называют
\emph{рядом Тейлора} функции $f$ с центром $x_0$ (независимо от того, сходится ли он к $f$).
\end{definition}

\section{Критерий аналитичности}

Запишем формулу Тейлора с остатком $R_n(x)=f(x)-\sum_{k=0}^n\frac{f^{(k)}(x_0)}{k!}(x-x_0)^k$.
Функция аналитична в $x_0$ тогда и только тогда, когда $R_n(x)\to0$ в окрестности $x_0$.

\begin{theorem}[достаточное условие]
Пусть $f\in C^\infty(x_0-\rho,x_0+\rho)$ и существуют $C>0$, $R>0$ такие, что для всех $n$ и всех
$x$ из окрестности
\[
  \abs{f^{(n)}(x)}\le C\,\frac{n!}{R^n}.
\]
Тогда $f$ аналитична на $(x_0-R,x_0+R)$.
\end{theorem}

\begin{proof}
Остаток в форме Лагранжа: при некотором $\xi$ между $x_0$ и $x$
\[
  \abs{R_n(x)}=\abs[\Big]{\frac{f^{(n+1)}(\xi)}{(n+1)!}(x-x_0)^{n+1}}
   \le C\frac{(n+1)!}{R^{n+1}}\cdot\frac{\abs{x-x_0}^{n+1}}{(n+1)!}
   =C\abs[\Big]{\frac{x-x_0}{R}}^{n+1}.
\]
Если $\abs{x-x_0}<R$, то $\abs{(x-x_0)/R}<1$ и $R_n(x)\to0$, т.е. ряд Тейлора сходится к $f$.
\end{proof}

\begin{remark}
Достаточно и более простой оценки $\abs{f^{(n)}(x)}\le C\,M^{n}$: тогда
$\abs{R_n}\le C\,M^{n+1}\abs{x-x_0}^{n+1}/(n+1)!\to0$ при любом $x$ (факториал перебивает
степень). Так доказывается аналитичность $e^x$, $\sin x$, $\cos x$ на всей оси.
\end{remark}

\section{Гладкость без аналитичности: функция Коши}

\begin{example}
\[
  f(x)=\begin{cases}e^{-1/x^2},& x\ne0,\\ 0,& x=0,\end{cases}
\]
принадлежит $C^\infty(\R)$, но \emph{не} аналитична в нуле. Все производные в нуле равны нулю:
$f^{(k)}(0)=0$ (каждая производная имеет вид $P(1/x)e^{-1/x^2}$, а $e^{-1/x^2}\to0$ быстрее любой
степени). Поэтому ряд Тейлора в нуле тождественно равен $0$ и не совпадает с $f$ ни в одной
проколотой окрестности. Это показывает, что $C^\omega\subsetneq C^\infty$.
\end{example}

\section{Теорема единственности}

\begin{theorem}[принцип изолированных нулей / единственности]
Пусть $f,g$ аналитичны на интервале $I$ и множество $E=\{x\in I: f(x)=g(x)\}$ имеет хотя бы одну
предельную точку в $I$. Тогда $f\equiv g$ на всём $I$.
\end{theorem}

\begin{proof}
Положим $h=f-g$ (аналитична); покажем $h\equiv0$. Пусть $x_0\in I$ — предельная точка $E$, т.е.
$h(x_0)=0$ и $x_0$ — предел нулей $h$. Если $h\not\equiv0$ около $x_0$, то
$h(x)=\sum_{n\ge k}a_n(x-x_0)^n=(x-x_0)^k\varphi(x)$, где $a_k\ne0$ — первый ненулевой
коэффициент и $\varphi(x_0)=a_k\ne0$. По непрерывности $\varphi\ne0$ в некоторой окрестности
$x_0$, значит $h(x)\ne0$ при $0<\abs{x-x_0}<\delta$ — у $x_0$ нет других нулей, противоречие с
тем, что $x_0$ предельная для $E$. Следовательно, все $a_n=0$, $h\equiv0$ в окрестности $x_0$.
Множество точек, где $h\equiv0$ вместе с окрестностью, одновременно открыто и замкнуто в $I$,
поэтому совпадает со всем $I$.
\end{proof}

\begin{remark}[аналитическое продолжение]
Из теоремы единственности следует: аналитическая функция полностью определяется своими значениями
на любом множестве с предельной точкой. Это даёт идею \emph{аналитического продолжения} — функцию,
заданную рядом в одной области, можно однозначно «продолжить» в большую область, и результат не
зависит от пути продолжения.
\end{remark}
